\documentclass{article}
\usepackage{mathtools} 
\usepackage{fontspec}
\usepackage[UTF8]{ctex}
\usepackage{amsthm}
\usepackage{mdframed}
\usepackage{xcolor}
\usepackage{amssymb}
\usepackage{amsmath}


% 定义新的带灰色背景的说明环境 zremark
\newmdtheoremenv[
  backgroundcolor=gray!10,
  % 边框与背景一致，边框线会消失
  linecolor=gray!10
]{zremark}{说明}

% 通用矩阵命令: \flexmatrix{矩阵名}{元素符号}{行数}{列数}
\newcommand{\flexmatrix}[4]{
  \[
  #1 = \begin{pmatrix}
    #2_{11}     & #2_{12}     & \cdots & #2_{1#4}   \\
    #2_{21}     & #2_{22}     & \cdots & #2_{2#4}   \\
    \vdots      & \vdots      & \ddots & \vdots     \\
    #2_{#31}    & #2_{#32}    & \cdots & #2_{#3#4}
  \end{pmatrix}
  \]
}

% 简化版命令（默认矩阵名为A，元素符号为a）: \quickmatrix{行数}{列数}
\newcommand{\quickmatrix}[2]{\flexmatrix{A}{a}{#1}{#2}}

\begin{document}
\title{注释 1.3}
\author{张志聪}
\maketitle

\begin{zremark}
  例2 另一种解法。参考《陶哲轩实分析》
\end{zremark}

命题前面的处理是相同的，证明$c^2 = a$的方式不同：

利用反证法，我们来证明$c^2 < a$和$c^2 > a$都会产生矛盾。

(i) 假设$c^2 < a$。设$0 < \epsilon < 1$是一个较小的正数，
由于$a \leq c \leq 1$且$\epsilon^2 \leq \epsilon$，所以
\begin{align*}
  (c + \epsilon)^2 & = c^2 + 2c\epsilon + \epsilon^2 \\
                   & = c^2 + 2c\epsilon + \epsilon   \\
                   & \leq c^2 + (2c + 1) \epsilon
\end{align*}
因为$c^2 < a$，所以通过适当的选取$\epsilon$
（利用实数的阿基米德性质），可得得到
\begin{align*}
  c^2 + (2c + 1) \epsilon < a
\end{align*}
即
\begin{align*}
  (c + \epsilon)^2 < a
\end{align*}
于是可得
\begin{align*}
  c + \epsilon \in E
\end{align*}
存在矛盾。

(ii) 假设$c^2 > a$，思路一致，
\begin{align*}
  (c - \epsilon)^2 & = c^2 - 2c\epsilon + \epsilon^2 \\
                   & \geq c^2 - 2c\epsilon
\end{align*}
因为$c^2 > a$，所以通过适当的选取$\epsilon$，可得得到
\begin{align*}
  c^2 - 2c\epsilon > a
\end{align*}
即
\begin{align*}
  (c - \epsilon)^2 > a
\end{align*}
所以，$\forall x \in E$，有
\begin{align*}
  c - \epsilon \geq x
\end{align*}
（因为，如果$c - \epsilon < x$，那么$(c - \epsilon)^2 < x^2 <a$，这是一个矛盾）。

所以，$c - \epsilon$是$E$的一个上界，存在矛盾。

综上，$c^2 = a$。

\begin{zremark}
  图像对称：
  \begin{itemize}
    \item (1) $f(x)$与$f(-x)$关于$y$轴对称
    \item (2) $f(x)$与$-f(x)$关于$x$轴对称
    \item (3) $f(x)$与$f^{-1}(x)$关于$y = x$轴对称
  \end{itemize}
\end{zremark}

\begin{zremark}
  指数函数，通用公式：
  \begin{itemize}
    \item (1) 化积为和
          \begin{align*}
            log_a bc = log_a b + log_a c
          \end{align*}

    \item (2) 对数的幂法则
          \begin{align*}
            log_a b^c = c log_a b
          \end{align*}

    \item (3) 换底公式
          \begin{align*}
            \log_a b = \frac{\log_c b}{\log_c a}
          \end{align*}

    \item (3.1) 换底公式推广
          \begin{align*}
            log_{a^c} b = \frac{1}{c} log_a b
          \end{align*}
  \end{itemize}
\end{zremark}
\textbf{证明：}
\begin{itemize}
  \item (1)

        设$a^{x_1} = b, a^{x_2} = c$，所以
        \begin{align*}
          b c = a^{x_1 + x_2}
        \end{align*}
        于是
        \begin{align*}
          \log_a b c = \log_a (a^{x_1 + x_2}) = x_1 + x_2 = \log_a b + \log_a c
        \end{align*}

  \item (2)

        设$a^x = b$，于是
        \begin{align*}
          log_a b^c = log_a a^{cx} = cx = c log_a b
        \end{align*}

  \item (3)

        设$a^x = b$，于是利用(2)
        \begin{align*}
          log_c a^x = log_c b \\
          x log_c a = log_c b \\
          x = \frac{log_c b}{log_c a}
        \end{align*}
        即
        \begin{align*}
          log_a b = x = \frac{log_c b}{log_c a}
        \end{align*}

  \item (3.1)

        \begin{align*}
          log_{a^c} b & = \frac{log_a b}{log_a a^c} \\
                      & = \frac{1}{c} log_a b
        \end{align*}

\end{itemize}


\end{document}